%% ---------------------------------------------------------------------
\section{Especificaciones de error y de transitorio}
\label{sec:er-diseno}
%% ---------------------------------------------------------------------

Una especificación de error se traduce en una cota sobre una constante: ``$\ess\le0.05$ ante rampa unitaria'' es $\kv\ge20$. Con un controlador proporcional, la única perilla es $K$, y la misma perilla mueve los polos del lazo cerrado. Las dos especificaciones pueden chocar.

\begin{ejemplo}[label=ej:er-conflicto]{Un servo que no puede con las dos}
Con $G_p=\dfrac{1}{s(s+4)(s+8)}$ y $\Gc=K$, se pide $\ess\le0.05$ ante rampa unitaria y $\OS\le20$\,\% ante escalón.
\solucion Tipo 1 con $\kv=K/32$: el error pide $K\ge640$. Pero la ecuación característica $s^3+12s^2+32s+K=0$ exige $K<12\cdot32=384$: \alerta{con $K$ solo, ni siquiera es estable}. Y el sobrepico pide un $K$ mucho menor: para $\zeta=0.49$ en los polos dominantes ($17$\,\% estimado, con margen; la simulación da $16$\,\%), $K=67.8$, que da $\kv=2.1$ y $\ess=0.47$ (\cref{fig:er-atraso}, azul).

Un \textbf{compensador de atraso} $\Gc=K\dfrac{s+a}{s+b}$ con $a>b>0$ multiplica $\kv$ por $a/b$. Si el par polo-cero está muy cerca del origen, en la zona de los polos dominantes $(s+a)/(s+b)\approx1$, y el transitorio casi no cambia. Se escoge $K=67.8$ por el transitorio, y $a/b=1.1\cdot640/67.8=10.4$ (con un $10$\,\% de margen); con $b=0.005$, $a=0.0519$. El lazo compensado tiene $\kv=22$ y $\ess=0.045$, con $\OS=19$\,\%; el precio es una cola lenta (un polo en $-0.053$ casi cancelado por el cero en $-0.052$): el tiempo de establecimiento sube de $3.1$ a $4.8$\,s, y el error de rampa tarda unos $116$\,s en quedar a menos del $2$\,\% de su valor final.
\end{ejemplo}

\begin{figure}[ht]
\centering
\includegraphics[width=\linewidth]{cap06/f_atraso}
\caption{Ejemplo~\ref{ej:er-conflicto}: con solo $K=67.8$ (azul) y con el atraso (coral). El atraso deja casi igual la respuesta al escalón y baja el error de rampa de $0.47$ a $0.045$, con una cola lenta.}
\label{fig:er-atraso}
\end{figure}

\codigorepo{cap06_error/C6_Ej5_Atraso.m}{conflicto}{el conflicto entre las dos especificaciones}
\codigorepo{cap06_error/C6_Ej5_Atraso.m}{atraso}{el compensador de atraso}

El diseño de atrasos en el dominio de la frecuencia es tema de la \cref{sec:rf-atraso}; el PI, que lleva $b$ hasta cero, del \cref{cap:pn}.
